\chapter*{\hfill{\centering  ЗАКЛЮЧЕНИЕ}\hfill}
\addcontentsline{toc}{chapter}{ЗАКЛЮЧЕНИЕ}

Исследование показало, что время выполнения алгоритмов Левенштейна и Дамерау~---~Левенштейна растёт геометрически с увеличением длины строк. Нерекурсивная реализация и рекурсивная с кэшем для алгоритма Левенштейна показывают лучшее время.

Цель лабораторной работы была достигнута: исследованы алгоритмы поиска расстояний Левенштейна и Дамерау~---~Левенштейна.

Выполнены следующие задачи:
\begin{enumerate}
    \item описать алгоритмы поиска расстояний Левенштейна и Дамерау~---~Левенштейна для нахождения редакционного расстояния между строками;
    \item реализовать данные алгоритмы;
    \item выполнить сравнительный анализ алгоритмов по затрачиваемым ресурсам (времени);
    \item описать и обосновать полученные результаты в отчете.
\end{enumerate}